% Generated from validated Article Draft Result. Do not hand-edit; revise the structured draft instead.

2022年のmultimodal化を一つのmodel familyとして扱うことはできない。Flamingoは画像と言語をまたぐvisual-language modelingの流れを代表し、DALL-E 2やImagenはtext-to-image generationの流れに位置する。前者は複数modalitiesを入力・理解・応答へ結び付ける問題、後者は条件から新しいmediaを生成する問題であり、技術的な目的も公開形態も異なる。『画像を扱う』という共通点だけで同列に並べると、この分岐を見失う。\autocite{src-13aed13ff36c,src-16a1a1dc486e,src-27297a80fac0,src-2c38d2f354fa,src-86ce6205a4ed,src-929080260e16,src-c5c7eddb6cf6,src-cf669aed6653}


生成メディアでは、model qualityだけでなく『どの形で触れられるか』が年内に大きく分岐した。DALL-E 2は研究発表からbeta、より広い利用、APIへとlifecycleを進め、Stable Diffusionは公開weightを通じてlocal executionや派生利用へつながる別のdistribution surfaceを形成した。研究発表、beta、API、open-weight releaseは同じ公開状態ではないため、本版では系列として関連付けつつもchronology上は別イベントとして扱う。\autocite{src-16a1a1dc486e,src-26aadbc0393f,src-27297a80fac0,src-2832ee43ab8f,src-2c38d2f354fa,src-86ce6205a4ed,src-929080260e16}


model releaseの外側では、dataとpersonalizationが生態系の独立要素として見えるようになった。LAION-5Bは大規模なimage-text data資産の存在を前景化し、DreamBoothは既存のtext-to-image modelを特定subjectへ適応させるpersonalization方向を示した。生成品質の比較だけではなく、どのdataを使えるか、既存modelをどこまで個別化できるかが、生成メディアの利用可能性を左右する層になった。\autocite{src-26aadbc0393f,src-2832ee43ab8f,src-76875f66d5f3,src-e0a494dc6724}


同じ年、生成対象は静止画像だけに留まらなかった。AudioLMは音声生成、Make-A-Videoはvideo generation、DreamFusionはtextから3D表現を得る方向をそれぞれ示した。これらを『次世代Stable Diffusion』のような一本の系譜にまとめるのは不適切である。modalityごとに表現形式、学習data、時間的一貫性や3D整合性など異なる制約があり、2022年はむしろ生成メディアの設計空間が複数方向へ枝分かれした年として読める。\autocite{src-1b692f6d18b5,src-b2dbb3098853,src-c19139f0be76}


この変化をまとめる鍵は『一番良い画像modelはどれか』ではなく、understanding、generation、data、personalization、distribution、modality expansionという層が分離したことである。特にopen-weight distributionは、frontier modelの規模とは別に、研究者・開発者・個人がmodelを実行・改変するsurfaceを広げた。2022年後半の生成メディアは、単発の研究demoからmodelとdataと利用形態が相互作用する生態系へ移り始めた。\autocite{src-13aed13ff36c,src-16a1a1dc486e,src-1b692f6d18b5,src-26aadbc0393f,src-27297a80fac0,src-2832ee43ab8f,src-2c38d2f354fa,src-76875f66d5f3,src-86ce6205a4ed,src-929080260e16,src-b2dbb3098853,src-c19139f0be76,src-c5c7eddb6cf6,src-cf669aed6653,src-e0a494dc6724}


\begin{claimboundary}
各systemの画質、能力、benchmark、一般化性能は著者・project・vendorの一次資料に帰属する。本稿は異なるmodalityや公開条件の結果を横断して独立再現した優劣へ変換せず、2022年に現れた設計層とdistribution境界を比較する。\autocite{src-16a1a1dc486e,src-26aadbc0393f,src-27297a80fac0,src-2832ee43ab8f,src-2c38d2f354fa,src-86ce6205a4ed,src-929080260e16}
\end{claimboundary}
